\section{Computations used in the proofs}\label{app:verification}

This appendix lists the computations on which the proofs depend. Statements proved by hand, in particular those of \cref{sec:theta,sec:boundary} and the evaluations in \cref{thm:r2}, use no computation. The data and the programs that carry out the computations below are available in the repository cited under Data availability.

\emph{Cusp polynomials.} In every rank, the obstruction polynomials are obtained from the procedure of \cref{sec:cusp} in exact arithmetic over $K=\Q(t)$, $t=1-\cos(\pi/(r+2))$. The saddle point is obtained from $Q_r=t$ by solving the saddle equations backwards, and these equations are then checked exactly in $K$. In rank nine, the polynomials are computed modulo $p$ at the five roots of $f$, as described in the proof of \cref{thm:r9}.

\emph{Ranks three, four and five.} The eliminations of \cref{lem:r3elim,lem:r4elim,lem:r5elim} are proved by rational ideal-membership certificates and Gr\"obner bases. Each identity is checked by exact rational expansion, and the reverse containments by exact polynomial division.

\emph{Ranks two, six, seven and nine.} These eliminations use \cref{lem:transfer}. For each rank, two computations are needed over $\F_p$: the quotient dimension of the affine obstruction ideal, and the finiteness of the quotient by the ideal of highest homogeneous parts. The primes are $p=1000003$ in ranks two, six and seven, and $p=1000033$ in rank nine. In rank nine, the local integrality at $p$ and the preservation of degrees under reduction are part of the proof of \cref{thm:r9}. In ranks two and six, the rational zeros are located on a finite grid and then verified by exact substitution; since their number equals the bound from \cref{lem:transfer}, no other zeros exist.

\emph{Companion certificates.} \Cref{lem:certred} relies on the certificate coefficients for $k=3,\dots,6$, and identity \eqref{eq:certid} is checked by exact expansion in the Laurent polynomial ring over $\Z$. The fifth and sixth companions use separate certificates at $u=1$ and $u=w$. For $k=6$ the certificate has been cleared of an explicitly nonzero polynomial denominator $D_{6,v}(w)$, and the identity theorem justifies cancellation across its zeros. \Cref{thm:evencomp} uses the certificates at $u=1$. The certificate for $k=2$ is printed in full in \eqref{eq:k2cert}.
